% Осень26!A9; нумерация внутри столбца без пустых строк.
% Источник: лекция 3, §14 (строка 378)
% Источник: лекция 3, §15 (строка 404)
\ticket{07}{Теорема об обратимости линейного отображения, близкого к обратимому}
\label{ticket:07}

\subsection*{\texorpdfstring{Множество обратимых операторов \(\Omega_m\); лемма об оценке снизу}{Множество обратимых операторов \textbackslash Omega\_m; лемма об оценке снизу}}

\begin{definition}

\(\Omega_m:=\{L\in\mathrm{Lin}(\mathbb{R}^m,\mathbb{R}^m):\ \exists L^{-1}\}\) --- множество обратимых линейных операторов из \(\mathbb{R}^m\) в \(\mathbb{R}^m\). Обратимость равносильна \(\det L\ne0\) (формулы Крамера / присоединённая матрица дают обратный оператор, который снова линеен и лежит в том же пространстве --- «а как ещё»).

Пространство операторов «сверхмногомерное»: операторы задаются матрицами \(m\times m\), т.е. живут в \(\mathbb{R}^{m^2}\), --- «чисто психологическая история».

\end{definition}

\begin{lemma}

Пусть \(B\in\mathrm{Lin}(\mathbb{R}^m,\mathbb{R}^m)\) и нашлась константа \(c>0\) такая, что \[\forall x:\qquad |Bx|\ge c\,|x|.\] Тогда \(B\in\Omega_m\) (обратим), и более того \[\|B^{-1}\|\le\frac1c .\]

Смысл: обычно норму оценивали сверху, а здесь оценка \textbf{снизу} --- «оператор не может достаточно крупный вектор отправить почти в ноль: если вектор длины 1, то уж длины \(c\) образ получится». «Основное мастерство --- это сформулировать; доказательство очевидно».

\end{lemma}

\begin{proof}

\begin{enumerate}
\def\labelenumi{\arabic{enumi}.}
\tightlist
\item
  \emph{Ядро тривиально.} Если \(Bx=0\), то \(0=|Bx|\ge c|x|\), значит \(x=0\): \(\ker B=\{0\}\).
\item
  \emph{Обратимость.} Для оператора \(\mathbb{R}^m\to\mathbb{R}^m\) тривиальность ядра, \(\det\ne0\) и биективность --- «пересказ одной и той же истории»; «магия квадратных матриц»: из инъективности следует сюръективность. Значит, \(B\) --- биекция \(\mathbb{R}^m\) на \(\mathbb{R}^m\), существует обратный линейный оператор \(B^{-1}\).
\item
  \emph{Оценка нормы.} Перепишем данное неравенство в терминах \(y:=Bx\): раз \(B\) --- биекция, при переборе всех \(x\) величины \(y=Bx\) пробегают всё \(\mathbb{R}^m\), а \(x=B^{-1}y\). Тогда для всех \(y\) \[|y|\ge c\,|B^{-1}y|\qquad\Longleftrightarrow\qquad |B^{-1}y|\le\frac1c\,|y|,\] а оценка \(|B^{-1}y|\le C|y|\) при всех \(y\) --- это в точности \(\|B^{-1}\|\le C\) (обсуждалось при изучении норм).
\end{enumerate}

\end{proof}

\textbf{Замечание (оценка снизу для обратимого).} Если \(A\in\Omega_m\), то \[\forall x:\qquad |Ax|\ge\frac{1}{\|A^{-1}\|}\,|x|.\] Доказательство: \(|x|=|A^{-1}(Ax)|\le\|A^{-1}\|\cdot|Ax|\), откуда \(|Ax|\ge|x|/\|A^{-1}\|\). «Как добывать такие оценки в обратную сторону? Всё вертится в рамках практически тех же формул».

\subsection*{Теорема об обратимости оператора, близкого к обратимому}

«Теорема с таким загадочным названием: об обратимости оператора, близкого к обратимому. Лучше уж такое запомнить, чем то, что в самой лемме».

\begin{theorem}

Пусть \(L\in\Omega_m\) и \(M\in\mathrm{Lin}(\mathbb{R}^m,\mathbb{R}^m)\) «близок» к \(L\): \[\|L-M\|<\frac{1}{\|L^{-1}\|}.\] (Близость мерим расстоянием в пространстве операторов --- нормой разности; константа \(1/\|L^{-1}\|\) --- «удачная»: в одномерной модели \(x\mapsto ax\) это ровно та константа, с которой всё проходит; радиус окрестности зависит от точки \(L\).) Тогда:

\begin{enumerate}
\def\labelenumi{\arabic{enumi}.}
\tightlist
\item
  \(M\) обратим: \(M\in\Omega_m\);
\item
  \(\|M^{-1}\|\le\dfrac{1}{\|L^{-1}\|^{-1}-\|L-M\|}\);
\item
  \(\|L^{-1}-M^{-1}\|\le\dfrac{\|L^{-1}\|}{\|L^{-1}\|^{-1}-\|L-M\|}\cdot\|L-M\|\)
\end{enumerate}

(п. 3: «насколько далеко друг от друга обратные операторы --- примерно расстояние между исходными, умноженное на техническую константу, зависящую только от \(L\); при малом \(\|L-M\|\) вылезает фактически \(\|L^{-1}\|\) в квадрате; в малой окрестности расстояние между обратными пропорционально расстоянию между исходными»).

\end{theorem}

\subsubsection*{Доказательство пунктов 1--2}

Целимся в лемму §14: оценим \(|Mx|\) снизу. Расщепим: \[Mx=Lx+(M-L)x,\] откуда «неравенство треугольника в отрицательную сторону» (большое минус маленькое --- «по смыслу \(\|M-L\|\) мало»): \[|Mx|\ge|Lx|-|(M-L)x|.\] Оценим слагаемые: по замечанию из §14 \(|Lx|\ge\dfrac{1}{\|L^{-1}\|}|x|\) («занизили первое слагаемое по замечанию к лемме»); по определению нормы оператора \(|(M-L)x|\le\|M-L\|\cdot|x|\). Итак, \[|Mx|\ge\Bigl(\frac{1}{\|L^{-1}\|}-\|L-M\|\Bigr)|x|=:c\,|x|,\] и \(c>0\) по условию теоремы («в скобке --- разность частей того неравенства»). По лемме §14 оператор \(M\) обратим и \(\|M^{-1}\|\le 1/c=\dfrac{1}{\|L^{-1}\|^{-1}-\|L-M\|}\). \(\blacksquare\)

\subsubsection*{Доказательство пункта 3}

Ключевое тождество; вещественный прототип --- формула \(\dfrac1M-\dfrac1L=\dfrac{L-M}{ML}\): \[M^{-1}-L^{-1}=M^{-1}\,(L-M)\,L^{-1}.\] Проверка --- «раскройте скобочку»: \(M^{-1}LL^{-1}-M^{-1}ML^{-1}=M^{-1}-L^{-1}\). «Операторы образуют кольцо: есть сложение и умножение матриц, прямо как у вещественных чисел; только умножение некоммутативное, и обратимость хромает --- полно матриц без обратных, в отличие от чисел, где необратим лишь нуль».

Оценим норму с помощью мультипликативности (\(\|BA\|\le\|B\|\cdot\|A\|\), лекция 2) и пункта 2: \[\|M^{-1}-L^{-1}\|\le\|M^{-1}\|\cdot\|L-M\|\cdot\|L^{-1}\|\le\frac{\|L^{-1}\|\cdot\|L-M\|}{\|L^{-1}\|^{-1}-\|L-M\|}. \qquad\blacksquare\]
